\documentclass[12pt,aspectratio=169]{beamer}
\input{../../mybeamer}

\title{Class 10: Faraday's Law, Part 2}
\subtitle{AP Physics C: E\&M}
\input{../../me}
\input{../../mycommands}


\begin{document}

\begin{frame}
  \maketitle
\end{frame}


\section{Additional Notes on Motional EMF: Work and Energy}

\begin{frame}{Motional EMF}
  \begin{columns}
    \column{.38\textwidth}
    \begin{tikzpicture}
      \foreach \x in {1,...,6}
      \foreach \y in {0,...,3} \node[cyan] at (\x,\y) {$\bm\times$};
      \draw[line width=.6mm] (1.3,.5)--(6,.5);
      \draw[line width=.6mm] (1.3,2.5)--(6,2.5);
      \draw[line width=1.7mm,red!60!black] (4.5,.3)--(4.5,2.7);
      \draw[thick,<->] (2.8,.5)--(2.8,2.5) node[midway,fill=black!2]{$\ell$};
      \draw[thick,|<->] (1.3,.2)--(4.5,.2) node[pos=.42,fill=black!2]{$x$};
      \draw[vectors] (4,3.1)--(5,3.1) node[right]{$\vec v$};
      \draw[thick] (1.3,.47) to[R=$R$] (1.3,2.53);
      \node[cyan] at (1.5,3.2) {$\vec B_\text{in}$};
    \end{tikzpicture}
    
    \column{.62\textwidth}
    When the rod slides to the right with constant velocity $\vec v$, magnetic
    flux $\Phi_m$ through the circuit loop is:

    \eq{-.15in}{
      \Phi_m=B{\color{magenta}A}=B{\color{magenta}\ell x}
    }

    \vspace{-.17in}The induced emf inside the circuit is:
    
    \eq{-.1in}{
      |\mathcal E|
      =\left|\diff{\Phi_m}t\right|
      =B\ell{\color{violet}\diff xt}
      =B\ell{\color{violet}v}
    }

    The induced current through the circuit can be calculated using Ohm's law:

    \eq{-.15in}{
      I=\frac{\mathcal E}R=\frac{B\ell v}R
    }
  \end{columns}
\end{frame}



\begin{frame}{Induced Current from Motional EMF}
  Induced current $I$ runs counter-clockwise, because (by Lenz's law) it must
  produce a magnetic field that \emph{opposes the increase in magnetic flux}.

  \vspace{.1in}
  \begin{columns}
    \column{.38\textwidth}
    \begin{tikzpicture}
      \foreach \x in {1,...,6}
      \foreach \y in {0,...,3} \node[cyan] at (\x,\y) {$\bm\times$};
      \draw[lightgray,line width=.6mm] (1.3,.5)--(6,.5);
      \draw[lightgray,line width=.6mm] (1.3,2.5)--(6,2.5);
      \draw[line width=1.7mm,lightgray] (4.5,.3)--(4.5,2.7);
      \draw[vectors] (4,3.1)--(5,3.1) node[right]{$\vec v$};
      \draw[thick,gray] (1.3,.47) to[R=$R$] (1.3,2.53);
      \node[cyan] at (1.5,3.2) {$\vec B_\text{in}$};
      \draw[violet,vectors] (4.5,.7)--(4.5,2.2) node[above]{$I$};
      \draw[violet,vectors] (4,2.5)--(2,2.5) node[midway,above]{$I$};
      \draw[violet,vectors] (2,.5)--(4,.5) node[midway,above]{$I$};
    \end{tikzpicture}

    \column{.62\textwidth}
    The power dissipated by the (ohmic) resistor is:

    \eq{-.1in}{
      P=I^2R=\left[\frac{B\ell v}R\right]^2R=\frac{B^2\ell^2v^2}R
    }

    \textbf{But where did the energy come from?}
  \end{columns}
\end{frame}



\begin{frame}{Magnetic Force on Induced Current}
  \begin{columns}
    \column{.38\textwidth}
    \begin{tikzpicture}
      \foreach \x in {1,...,6}
      \foreach \y in {0,...,3} \node[cyan] at (\x,\y) {$\bm\times$};
      \draw[line width=.6mm] (1.3,.5)--(6,.5);
      \draw[line width=.6mm] (1.3,2.5)--(6,2.5);
      \draw[line width=1.7mm,red!70!black] (4.5,.3)--(4.5,2.7);
      \draw[thick,<->] (2.8,.5)--(2.8,2.5) node[midway,fill=black!2]{$\ell$};
      \draw[vectors] (4,3.1)--(5,3.1) node[right]{$\vec v$};
      \draw[vectors,orange] (4.58,1.5)--(5.5,1.5) node[right]{$\vec F_a$};
      \draw[thick] (1.3,.47) to[R=$R$] (1.3,2.53);
      \node[cyan] at (1.5,3.2) {$\vec B_\text{in}$};
      \foreach \yy in {.6,.9,...,2.5}
      \draw[orange,vectors] (4.42,\yy)--(3.8,\yy);
      \node[left,orange] at (3.9,1.5){$\vec F_m$};
      \draw[black!10,vectors] (4.5,.7)--(4.5,2.2) node[above]{$I$};
    \end{tikzpicture}
    
    \column{.62\textwidth}
    As the induced current $I$ in the rod moves in the magnetic field $\vec B$,
    the rod must experience a magnetic force $\vec F_m$ with magnitude:

    \eq{-.15in}{
      F_m=I\ell B = \left[\frac{B\ell v}R\right]\ell B
      =\frac{B^2\ell^2v}R
    }

    For a constant $\vec v$, an external force of $\vec F_a=-\vec F_m$ must be
    applied (1st law of motion). The power generated by the applied force is:

    \eq{-.2in}{
      P = F_av=F_mv=\left[\frac{B^2\ell^2v}R\right]v=\frac{B^2\ell^2v^2}R
    }
  \end{columns}
\end{frame}



\begin{frame}{Energy Conservation}
  \begin{columns}
    \column{.38\textwidth}
    \begin{tikzpicture}
      \foreach \x in {1,...,6}
      \foreach \y in {0,...,3} \node[cyan] at (\x,\y) {$\bm\times$};
      \draw[line width=.6mm] (1.3,.5)--(6,.5);
      \draw[line width=.6mm] (1.3,2.5)--(6,2.5);
      \draw[line width=1.7mm,red!70!black] (4.5,.3)--(4.5,2.7);
      \draw[vectors] (4,3.1)--(5,3.1) node[right]{$\vec v$};
      \draw[vectors,orange] (4.58,1.5)--(5.5,1.5) node[right]{$\vec F_a$};
      \draw[thick] (1.3,.47) to[R=$R$] (1.3,2.53);
      \node[cyan] at (1.5,3.2) {$\vec B_\text{in}$};
    \end{tikzpicture}
        
    \column{.62\textwidth}
    The power generated by the applied force is equal to the power dissipated
    by the resistor that we have calculated:

    \eq{-.1in}{
      P = \frac{B^2\ell^2v^2}R
    }

    The energy in the circuit comes from the positive work done by the applied
    force, and very importantly, our calculations confirm the law of
    conservation of energy.
  \end{columns}
\end{frame}



\section{Asymmetry in Faraday's Law}

\begin{frame}{A Coil Moving Through a Magnetic Field}
  A loop of wire moves through a magnetic field $\vec B$ into the page with
  \underline{constant} velocity $\vec v$. From Faraday's law, when the wire
  moves \emph{into} the magnetic field, and when it moves \emph{out} of the
  field, magnetic flux $\Phi_m$ changes, creating an electromotive force
  $\mathcal E$:

  \eq{-.1in}{
    \boxed{ \mathcal E=-\diff{\Phi_m}t }
  }
  \begin{center}
    \begin{tikzpicture}[scale=.6]
      \draw[cyan] (-.3,-.3) rectangle (8.3,5.3) node[right]{$\vec B_\text{in}$};
      \foreach \x in {0,1,...,8}{
        \foreach \y in {0,1,...,5} \node[cyan] at (\x,\y) {$\bm\times$};
      }
      
      \fill[gray,opacity=.1] (-.3,1.5) rectangle (1,3.5);
      \draw[very thick] (-1.5,1.5) rectangle (1,3.5);
      \draw[vectors] (1,2.5)--(2.2,2.5) node[above]{$\vec v$};
      \node[below] at (-.3,1.5) {$\Phi_m$ increases};

      \fill[gray,opacity=.1] (3,1.5) rectangle (5.5,3.5);
      \draw[very thick] (3,1.5) rectangle (5.5,3.5);
      \draw[vectors] (5.5,2.5)--(6.7,2.5) node[above]{$\vec v$};
      \node[below] at (4.25,1.5) {$\Phi_m$ constant};

      \fill[gray,opacity=.1] (7.5,1.5) rectangle (8.3,3.5);
      \draw[very thick] (7.5,1.5) rectangle (10,3.5);
      \draw[vectors] (10,2.5)--(11.2,2.5) node[above]{$\vec v$};
      \node[below] at (8.75,1.5) {$\Phi_m$ decreases};
    \end{tikzpicture}
  \end{center}
\end{frame}



\begin{frame}{Coil Moving Through a Magnetic Field}
  For an observer in the same frame of reference as the magnetic field, the
  reason is obvious. The charges\footnote{The charges are presumably positive,
  but it doesn't really matter} inside the wire experience a
  \emph{magnetic force} as they move through the magnetic field, creating an
  \emph{emf}, and an electric current. This is called \textbf{motional EMF}.
  \begin{center}
    \begin{tikzpicture}[scale=.65]
      \draw[cyan] (-.3,-.3) rectangle (8.3,5.3) node[right]{$\vec B_\text{in}$};
      \foreach \x in {0,1,...,8}{
        \foreach \y in {0,1,...,5} \node[cyan] at (\x,\y) {$\bm\times$};
      }
      
      \fill[gray,opacity=.1] (-.3,1.5) rectangle (1,3.5);
      \draw[very thick] (-1.5,1.5) rectangle (1,3.5);
      \draw[vectors] (1,2.5)--(2.2,2.5) node[above]{$\vec v$};
      \draw[vectors,blue] (-1.2,2) ..controls (1.5,1) and (1.5,4).. (-1.2,3)
      node[below]{$I$};

      \fill[gray,opacity=.1] (3,1.5) rectangle (5.5,3.5);
      \draw[very thick] (3,1.5) rectangle (5.5,3.5);
      \draw[vectors] (5.5,2.5)--(6.7,2.5) node[above]{$\vec v$};
      \node[below] at (4.25,1.5) {No current};

      \fill[gray,opacity=.1] (7.5,1.5) rectangle (8.3,3.5);
      \draw[very thick] (7.5,1.5) rectangle (10,3.5);
      \draw[vectors] (10,2.5)--(11.2,2.5) node[above]{$\vec v$};
      \draw[vectors,blue] (9.7,2) ..controls(7,1) and (7,4).. (9.7,3)
      node[below]{$I$};
    \end{tikzpicture}
  \end{center}
  \vspace{.2in}
\end{frame}




\begin{frame}{Magnetic Field Moving Past a Coil of Wire}
  For an observer in the same frame of reference as the wire, the reason is
  also obvious. As the magnetic field moves, the magnetic flux changes, and an
  \emph{electric field} is created in the wire. A current is created because
  the charges inside the wire experience an \emph{electric force}.
  \begin{center}
    \vspace{-.15in}
    \begin{tikzpicture}[scale=.7]
      \draw[cyan] (-.3,-.3) rectangle (8.3,5.3) node[right]{$\vec B_\text{in}$};
      \draw[vectors,cyan] (-.3,4)--(-1.5,4) node[left]{$\vec v$};
      \foreach \x in {0,1,...,8}{
        \foreach \y in {0,1,...,5} \node[cyan] at (\x,\y) {$\bm\times$};
      }
      
      \fill[gray,opacity=.1] (-.3,1.5) rectangle (1,3.5);
      \draw[very thick] (-1.5,1.5) rectangle (1,3.5);
      \draw[vectors,blue] (-1.2,2) ..controls (1.5,1) and (1.5,4).. (-1.2,3)
      node[below]{$I$};
    \end{tikzpicture}
  \end{center}
\end{frame}


\begin{frame}{Asymmetry in Faraday's Law}
  Both observers obtain the same result (same $\mathcal E$ and same current $I$
  in the wire), however they cannot agree on the \emph{reason}:
  \begin{itemize}
  \item In one inertial (non-accelerating) frame of reference, the \emph{emf}
    is due to magnetic force on the charges
  \item In another inertial frame, the \emph{emf} is due to an electric force
    on the same charges
  \end{itemize}
  Why?
\end{frame}


\begin{frame}{Asymmetry in Faraday's Law}
  \begin{itemize}
  \item As expected, results from applying Faraday's law is the same in both
    inertial frames of reference
    \begin{itemize}
    \item This is the application of the \textbf{principle of relativity} which
      states that all physical laws must be equivalent in all inertial frames
      of reference
    \item The principle is the foundation of all physical laws, and is known
      since Galileo's time
    \end{itemize}
  \item This suggests that this ``asymmetry'' in Faraday's law 
    \begin{itemize}
    \item is because magnetic force and electric force are not two distinct
      forces
    \item $\vec F_m$ is just $\vec F_q$ when viewed from a different frame of
      reference
    \end{itemize}
  \end{itemize}
\end{frame}



\begin{frame}{Special Relativity}
  Albert Einstein's 1905 paper
  \emph{On the electrodynamics of moving bodies}\footnote{\emph{Zur
  Elektrodynamik bewegter K\"{o}rper}, Annalen der Physik, 17(10):891--921,
  30 June, 1905}, his first paper on the theory of relativity, addresses this
  problem. In the first paragraph:
  \begin{center}
    \begin{minipage}{.8\textwidth}
      \footnotesize
      It is known that Maxwell's electrodynamics---as usually understood at the
      present time---when applied to moving bodies, leads to asymmetries which
      do not appear to be inherent in the phenomena\ldots if the magnet is
      in motion and the conductor at rest, there arises in the neighbourhood of
      the magnet and electric field with a certain finite energy, producing
      current at the places where parts of the conductor are situated. But if
      the magnet is stationary and and the conductor in motion, no electric
      field arises in the neighbourhood of the magnet. In the conductor,
      however, we find an electromotive force, to which itself there is no
      corresponding energy\ldots
    \end{minipage}
  \end{center}
  Only in subsequent sections where Einstein begins to mention the other
  aspects in special relativity that we usually teach the topic on.
  \vspace{.2in}
\end{frame}



\begin{frame}{Special Relativity}
  The bottom line is that electromagnetism can be explained
  \begin{itemize}
  \item Using electric force (coulomb's law), and
  \item Length contraction, which is one of the results of the relativity of
    space and time
  \end{itemize}
\end{frame}
\end{document}
